% !TEX program = pdflatex
% Example: arXiv preprint template
% Compile with: pdflatex main.tex && bibtex main && pdflatex main.tex && pdflatex main.tex
% Requires: arxiv style (included in TeX Live / MiKTeX)

\documentclass{article}

% Note: The arxiv.sty package (https://github.com/kourgeorge/arxiv-style)
% provides arXiv-like formatting. If not installed, this template
% compiles with standard article class formatting.

\usepackage{amsmath,amssymb}
\usepackage{graphicx}
\usepackage{booktabs}
\usepackage{url}

\title{A Framework for Self-Supervised Learning in Time Series Analysis}

\author{
  Author One\thanks{Corresponding author: \texttt{one@uni.edu}} \\
  Department of Computer Science, University A \\
  \texttt{one@uni.edu} \\
  \and
  Author Two \\
  Department of Computer Science, University B \\
  \texttt{two@uni.edu}
}

\begin{document}
\maketitle

\begin{abstract}
Self-supervised learning has emerged as a powerful paradigm
for learning representations from unlabeled data. We propose
a framework that adapts self-supervised techniques to time
series analysis, achieving state-of-the-art results on
multiple benchmarks.
\end{abstract}

\section{Introduction}
Self-supervised learning (SSL) has revolutionized natural
language processing and computer vision. However, its
application to time series data remains underexplored.

\section{Methodology}
Our framework consists of three components:
\begin{enumerate}
\item A pretext task based on temporal contrastive learning.
\item A transformer-based encoder architecture.
\item A fine-tuning strategy for downstream tasks.
\end{enumerate}

The contrastive loss is defined as:
\begin{equation}
\mathcal{L} = -\sum_{i} \log \frac{\exp(\mathbf{z}_i \cdot \mathbf{z}_j / \tau)}
{\sum_k \exp(\mathbf{z}_i \cdot \mathbf{z}_k / \tau)}
\end{equation}

\section{Experiments}
Table~\ref{tab:results} shows results on three benchmarks.

\begin{table}[htbp]
\caption{Classification accuracy on benchmark datasets.}
\label{tab:results}
\centering
\begin{tabular}{lccc}
\toprule
Method & Dataset A & Dataset B & Dataset C \\
\midrule
Supervised & 88.2 & 91.5 & 85.3 \\
Random Init & 72.1 & 78.3 & 70.8 \\
Proposed SSL & 90.7 & 93.2 & 88.9 \\
\bottomrule
\end{tabular}
\end{table}

\section{Conclusion}
We presented a self-supervised learning framework for time
series that outperforms supervised baselines while requiring
no labeled data for pre-training.

\bibliographystyle{plain}
\bibliography{references}

\end{document}
